\documentclass{article}

\usepackage[preprint]{neurips_2024}
\usepackage[utf8]{inputenc}
\usepackage[T1]{fontenc}
\usepackage{booktabs}
\usepackage{hyperref}
\usepackage{graphicx}
\usepackage{amsmath}

\title{Revisiting Token Merging for CogVideoX: Asymmetric Reduction and Restoration for Training-free Video Generation Acceleration}

\author{CS7352 Course Project Team}

\begin{document}
\maketitle

\begin{abstract}
We study training-free inference acceleration for CogVideoX-2B through visual Token Merging.
Early block-level hidden-state merging reduces computation but introduces blur, grid artifacts, flicker, and structural distortion.
We therefore implement a schedule-aware asymmetric reduction/restoration variant inspired by AsymRnR: visual Q and/or K/V tokens are reduced inside attention, Q outputs are restored to full visual length, and text tokens remain untouched.
This report skeleton intentionally leaves unrun metrics as ``not run yet'' until benchmark scripts are executed.
\end{abstract}

\section{Introduction}
Video diffusion transformers are expensive because attention is applied to long spatiotemporal token sequences.
The project follows the original proposal: accelerate open-source video generation models, primarily CogVideoX-2B, using training-free spatial and spatiotemporal Token Merging.
We retain inherited naive ToMe as a baseline and add attention-level asymmetric reduction/restoration to improve the quality-efficiency trade-off.

\section{Related Work}
\paragraph{Video diffusion transformers.}
CogVideoX uses a transformer-based video generation architecture with joint text/video attention.

\paragraph{Token Merging and ToMe.}
ToMe performs bipartite matching to merge similar tokens and can reduce attention cost without training.
Direct hidden-state merging, however, is risky in iterative video denoising.

\paragraph{Diffusion acceleration.}
Cache and distillation methods are important related work but are not the main implementation here.

\paragraph{AsymRnR and token reduction.}
AsymRnR motivates asymmetric token reduction and restoration in diffusion transformers.
Our implementation is a project-owned AsymRnR-lite variant for CogVideoX attention processors.

\section{Method}
\subsection{Naive Frame-wise Spatial ToMe Baseline}
The inherited baseline merges visual hidden states inside selected CogVideoX blocks and restores them after attention/FFN.
It is preserved as \texttt{--accel naive\_tome}.

\subsection{Problems with Hidden-state ToMe in Video DiT}
Block-level merging removes independent representations for absorbed visual positions and can amplify artifacts over denoising steps.
Observed artifacts include blur, pixelation, flicker, structure distortion, and motion instability.

\subsection{Asymmetric Q/V Reduction-Restoration}
The new method performs Q/K/V projection first, applies RoPE to full visual Q/K, reduces visual Q and/or K/V, runs SDPA, then restores Q-output to the original visual length.
Text tokens are never reduced.

\subsection{Spatiotemporal Bipartite Matching}
Visual tokens are partitioned by a destination stride over latent time, height, and width.
Source tokens are matched to destination tokens by cosine, dot, Euclidean, or random similarity.

\subsection{Euclidean Matching and Replace Reduction}
The default configuration uses Euclidean similarity and replace/discard reduction to avoid the extra averaging cost of mean reduction.
Mean reduction is retained as an ablation.

\subsection{Timestep, Block, and Feature Schedule}
Q and K/V can use different ratios.
Blocks are selected by layer strategy, and early/late denoising steps can be skipped.

\subsection{Matching Cache}
The runtime caches matching indices/plans by feature, block, shape, stride, ratio, and step window.
The cache is reset per prompt and stores no hidden features.

\section{Experiments}
\subsection{Model and Setting}
Model: CogVideoX-2B. Default formal setting: 49 frames, 480$\times$720, 50 denoising steps, seed 42.

\subsection{Prompts}
The benchmark prompt set is stored in \texttt{configs/prompts\_benchmark.txt}.

\subsection{Metrics}
Efficiency metrics: end-to-end latency, transformer latency, average denoising step time, peak GPU memory allocated/reserved, token reduction ratio, matching time, and cache hit rate.
Quality metrics: lightweight frame MSE/PSNR against same-seed baseline, temporal difference proxy, and human pairwise review.
FVD/VBench are not reported unless actually run.

\subsection{Baselines}
Baselines include no acceleration, inherited naive ToMe at multiple ratios, KV-only RnR, Q/V RnR, and final RnR-ToMe.

\begin{table}[t]
\centering
\caption{Main latency-quality comparison.}
\begin{tabular}{lrrrrr}
\toprule
Method & E2E s & Step s & Peak GiB & Quality metric & Status \\
\midrule
none & not run yet & not run yet & not run yet & not run yet & TODO \\
naive\_tome r=0.2 & not run yet & not run yet & not run yet & not run yet & TODO \\
kv\_rnr & not run yet & not run yet & not run yet & not run yet & TODO \\
qv\_rnr & not run yet & not run yet & not run yet & not run yet & TODO \\
rnr\_tome & not run yet & not run yet & not run yet & not run yet & TODO \\
\bottomrule
\end{tabular}
\end{table}

\begin{table}[t]
\centering
\caption{Ablation on similarity, reduction, and cache.}
\begin{tabular}{llllr}
\toprule
Similarity & Reduction & Cache steps & Variant & Status \\
\midrule
cosine & replace & 5 & rnr\_tome & not run yet \\
euclidean & replace & 5 & rnr\_tome & not run yet \\
euclidean & mean & 5 & rnr\_tome & not run yet \\
euclidean & replace & 1 & rnr\_tome & not run yet \\
\bottomrule
\end{tabular}
\end{table}

\begin{table}[t]
\centering
\caption{Prompt-category breakdown.}
\begin{tabular}{lrrrr}
\toprule
Category & Prompt count & Best speedup & Quality note & Status \\
\midrule
low motion & 2 & not run yet & TODO & TODO \\
high motion & 2 & not run yet & TODO & TODO \\
camera motion & 2 & not run yet & TODO & TODO \\
fine details & 2 & not run yet & TODO & TODO \\
multi-object & 2 & not run yet & TODO & TODO \\
\bottomrule
\end{tabular}
\end{table}

\begin{table}[t]
\centering
\caption{Memory and token statistics.}
\begin{tabular}{lrrrr}
\toprule
Method & Peak allocated GiB & Peak reserved GiB & Q reduction & KV reduction \\
\midrule
none & not run yet & not run yet & 0 & 0 \\
kv\_rnr & not run yet & not run yet & not run yet & not run yet \\
qv\_rnr & not run yet & not run yet & not run yet & not run yet \\
rnr\_tome & not run yet & not run yet & not run yet & not run yet \\
\bottomrule
\end{tabular}
\end{table}

\section{Qualitative Analysis}
Manual inspection should focus on blur, pixelation, flicker, structure distortion, and motion artifacts.
Pairwise review templates are generated by \texttt{scripts/evaluate\_quality.py}.

\section{Limitations}
The current RnR schedule is an AsymRnR-lite implementation with fixed ratios and explicit gates rather than a learned or pre-estimated similarity-threshold policy.
Full-resolution CogVideoX-2B is memory intensive on 12 GiB GPUs.

\section{Conclusion}
This work reframes the inherited naive ToMe baseline into a reproducible comparison point and adds attention-level asymmetric reduction/restoration for safer training-free acceleration.

\section{Contribution Statement}
Existing naive ToMe code was inherited from preliminary project work.
This implementation adds attention-level asymmetric RnR, scheduling, matching cache, unified benchmark scripts, evaluation scripts, audits, and report/result structure.
AsymRnR, Diffusers, and CogVideoX were used as references, not copied wholesale.

\bibliographystyle{plain}
\bibliography{references}

\end{document}
